\documentclass[11pt,a4paper]{article}

% ═══════════════════════════════════════════════════════════════════
%  PACKAGES
% ═══════════════════════════════════════════════════════════════════
\usepackage[utf8]{inputenc}
\usepackage[T1]{fontenc}
\usepackage{amsmath,amssymb,amsthm}
\usepackage{mathtools}
\usepackage{booktabs}
\usepackage{array}
\usepackage{graphicx}
\usepackage{xcolor}
\usepackage{hyperref}
\usepackage[margin=1in,headheight=14pt]{geometry}
\usepackage{enumitem}
\usepackage{fancyhdr}
\usepackage{titlesec}
\usepackage{algorithm}
\usepackage{algpseudocode}
\usepackage{listings}
\usepackage{caption}
\usepackage{subcaption}
\usepackage[numbers]{natbib}

% ═══════════════════════════════════════════════════════════════════
%  CONFIGURATION
% ═══════════════════════════════════════════════════════════════════
\hypersetup{
    colorlinks=true,
    linkcolor=blue!70!black,
    citecolor=green!50!black,
    urlcolor=blue!60!black
}

\lstset{
    basicstyle=\ttfamily\small,
    breaklines=true,
    frame=single,
    backgroundcolor=\color{gray!10},
    numbers=none,
    columns=fullflexible,
    keepspaces=true
}

\pagestyle{fancy}
\fancyhf{}
\rhead{Daugherty, Ward, Ryan \& Martin --- Beale Cipher Forensics}
\lhead{SmartLedger Research}
\rfoot{Page \thepage}

\newtheorem{theorem}{Theorem}
\newtheorem{lemma}[theorem]{Lemma}
\newtheorem{proposition}[theorem]{Proposition}
\newtheorem{corollary}[theorem]{Corollary}
\newtheorem{definition}{Definition}

% ═══════════════════════════════════════════════════════════════════
%  TITLE
% ═══════════════════════════════════════════════════════════════════
\title{%
    \textbf{Computational Forensic Analysis of the Beale Ciphers:}\\[6pt]
    \Large Spectral Diagnostics, Fabrication Scoring, Edition Reconstruction,\\
    and Multi-Channel Forensic Analysis via Ising Optimization
}

\author{
    Bryan Daugherty, CCI, CBI, SME\textsuperscript{1}
    \and
    Gregory Ward\textsuperscript{1}
    \and
    Shawn Ryan\textsuperscript{1}
    \and
    J. Alexander Martin\textsuperscript{1}
    \\[6pt]
    \textsuperscript{1}SmartLedger Research / Origin Neural AI
}

\date{March 2026}

\begin{document}
\maketitle

% ═══════════════════════════════════════════════════════════════════
%  ABSTRACT
% ═══════════════════════════════════════════════════════════════════
\begin{abstract}
We present a computational forensic analysis of the three Beale ciphers
(B1,~B2,~B3) spanning eight phases and 35+ statistical diagnostics.

\textbf{Methods.}
We employ a 4-solver Ising ensemble, Welch-averaged bispectral analysis,
QUBO-formulated edition search, parametric fabrication simulation with
bootstrap validation, and multi-channel forensic analysis extending the
trifurcating cipher methodology developed for Kryptos~K4~\cite{daugherty2026k4}.
Analysis code ($\sim$16{,}000 lines, 25~scripts) was executed on a
combinatorial optimization platform.

\textbf{Results.}
B2 is genuine (fabrication score~0.00, decode accuracy~96.6\%,
edition-corrected~$\sim$99\%+). B1 and~B3 are fabrications
(scores~4.01/7.51) with formal impossibility proofs
($\chi^2_{\text{B1}}{=}50.9$, $\chi^2_{\text{B3}}{=}57.8$, threshold~44.3).
Ward's construction timeline is reconstructed (B2$\to$B3$\to$B1,
$\sim$45~min total), cognitive identity profiling shows consistency
with a single fabricator ($p{=}0.052$, 12D dual-null bootstrap), all 22~B2 mismatches resolve into
4~systematic groups, and joint simulated annealing validates the
edition-specific offset model (92.2\% accuracy, 31\% offset rediscovery).
Four forensic information channels are identified: the genuine B2
surface encoding, a structural edition fingerprint, Ward's cognitive
fabrication signature, and the deliberately embedded Gillogly
alphabetical sequence.
\end{abstract}

% ═══════════════════════════════════════════════════════════════════
%  1. INTRODUCTION
% ═══════════════════════════════════════════════════════════════════
\section{Introduction}

The Beale ciphers, published in an 1885 pamphlet attributed to
James~B.~Ward~\cite{ward1885}, purportedly describe the location~(B1),
contents~(B2), and heirs~(B3) of a treasure buried in Bedford County,
Virginia in the 1820s. Only B2 has been decrypted, using the Declaration
of Independence (DoI) as a book cipher key, with first letters of
numbered words spelling out the plaintext.

Gillogly~\cite{gillogly1980} first noted the suspicious alphabetical run
in B1's DoI-decoded output. Fitzgerald~\cite{fitzgerald2024} introduced
quantitative fabrication metrics achieving Bayes Factor~$\sim 2 \times 10^7$.
Wase~\cite{wase2020} provided independent statistical analysis, and
Hammer~\cite{hammer1968} identified early structural signatures. However,
no prior analysis has combined spectral forensics, formal impossibility
proofs, cognitive identity testing, edition reconstruction, and
multi-channel forensic analysis into a unified framework.

We present such an analysis across eight phases comprising 35+~diagnostics
and $\sim$16{,}000 lines of analysis code, executed on a
combinatorial optimization platform. Our contributions include:

\begin{enumerate}[nosep]
    \item Formal impossibility proofs that B1/B3 encode gibberish under
          \emph{any} key text (Section~\ref{sec:impossibility})
    \item Identification of Beale's specific DoI edition via offset cluster
          analysis and joint SA optimization (Section~\ref{sec:edition})
    \item Cognitive identity profiling consistent with a single fabricator
          for B1 and~B3 (Section~\ref{sec:cognitive})
    \item Welch-averaged bicoherence correcting single-realization spectral
          saturation (Section~\ref{sec:spectral})
    \item Multi-channel forensic analysis applying the trifurcating cipher
          framework from K4~\cite{daugherty2026k4} (Section~\ref{sec:channels})
    \item Complete resolution of all 22~B2 mismatches into 4~systematic
          groups (Section~\ref{sec:edition})
\end{enumerate}

% ═══════════════════════════════════════════════════════════════════
%  2. RELATED WORK
% ═══════════════════════════════════════════════════════════════════
\section{Related Work}
\label{sec:related}

Gillogly~\cite{gillogly1980} identified the 12-letter alphabetical run
CDEFGHIIJKLM in B1's DoI-decoded output and noted that B1/B3 decoded
text has non-English letter frequencies. His analysis was qualitative;
we formalize this observation as a chi-squared impossibility proof
(Theorem~\ref{thm:impossibility}) and prove the Gillogly sequence was
deliberately planted via monotonicity analysis (Section~\ref{sec:gillogly}).

Fitzgerald~\cite{fitzgerald2024} introduced four quantitative fabrication
metrics (serial correlation, distinct ratio, fatigue gradient, page
boundary clustering), achieving Bayes Factor~$\sim 2 \times 10^7$ for
fabrication detection. We adopt and extend his metrics, adding changepoint
detection, chi-squared selection models, and a parametric 8-parameter
Ward construction simulator that reproduces B1/B3 diagnostics within
1--2$\sigma$ (Section~\ref{sec:wardmodel}).

Wase~\cite{wase2020} provided independent statistical analysis of number
distributions. Hammer~\cite{hammer1968} identified early structural
signatures in the cipher numbers.

Friedman's~\cite{friedman1987} index of coincidence provides one of our
three independent impossibility arguments. The Needleman-Wunsch
algorithm~\cite{needleman1970}, originally developed for protein sequence
alignment, serves as our primary tool for B2 decode accuracy measurement,
resolving the 796~vs~763 character discrepancy that misled earlier analyses
into reporting $\sim$34\% accuracy instead of the true 96.6\%.

Our prior work on Kryptos~K4~\cite{daugherty2026k4} established the
concept of a \emph{trifurcating cipher}---a single ciphertext encoding
multiple information channels through a shared invariant anchor. We extend
this framework to the Beale ciphers, distinguishing between intentional
and emergent information channels (Section~\ref{sec:channels}).

% ═══════════════════════════════════════════════════════════════════
%  3. METHODOLOGY
% ═══════════════════════════════════════════════════════════════════
\section{Methodology}
\label{sec:methodology}

Our analysis proceeds in eight phases, each motivated by open questions
from the preceding phase. The progressive structure is essential:
later phases build on calibrated baselines and resolved ambiguities
from earlier work.

\subsection{Phase~1: Statistical Fingerprinting}

We begin with twelve foundational diagnostics to establish baseline
distributional properties and identify anomalies warranting deeper
investigation. These include number distribution statistics,
Benford's Law testing, last-digit uniformity, Shannon entropy,
number frequency analysis, bigram analysis, autocorrelation
(50~lags), spectral fingerprinting (transition matrix eigenvalues
mod~26), index of coincidence~\cite{friedman1987}, Gillogly checks,
cross-cipher comparison, and page boundary analysis. The key
diagnostic outcome---B2's near-zero serial correlation (0.044) vs
B1/B3's elevated values (0.252/0.612)---motivates the fabrication
scoring framework in Phase~2.

\subsection{Phase~2: Fabrication Scoring}

Phase~1 anomalies are formalized using four Fitzgerald-class
metrics~\cite{fitzgerald2024} with Bayes Factor~$\sim 2 \times 10^7$.
We extend Fitzgerald's framework with a chi-squared selection model
that distinguishes \emph{letter-targeted} selection (genuine) from
\emph{DoI-uniform} selection (fabricated), two-phase segmentation with
changepoint detection revealing fatigue transitions, and Monte Carlo
null distributions ($N{=}5{,}000$) for calibration.

\subsection{Phase~3: Ward Construction Model}

To test whether a single fabricator model can reproduce B1/B3
diagnostics, we build an 8-parameter parametric simulator
(WardFabricationModel). If the model fits both ciphers, this supports
the single-fabricator hypothesis and constrains the construction
timeline. Gillogly forensics (monotonicity proof) and even-digit bias
investigation are conducted in this phase.

\subsection{Phase~4: Cryptanalysis \& Edition Identification}

We perform full B2 decryption with Needleman-Wunsch
alignment~\cite{needleman1970}, resolving the 796~vs~763 discrepancy.
DoI edition fingerprinting identifies 51~non-zero offsets revealing
$\sim$13 extra words in Beale's edition. B3~Q1 fragment recovery
rules out genuine content. Formal B1/B3 impossibility proofs are
established.

\subsection{Phase~5: Edition Reconstruction \& Mismatch Resolution}

All 22~true B2 mismatches are exhaustively resolved into 4~systematic
groups plus orphan errors. X-word and Y-word mysteries are investigated.
Ward cognitive forensics (sliding window cognitive load model) and
Gillogly construction simulation are performed.

\subsection{Phase~6: Spectral Forensics \& Publication Preparation}

A critical Y/X-word alignment bug is fixed (8/9 not 2/9 for Y-word).
Eight bispectral diagnostic tools are applied. QUBO edition search
with the 4-solver Ising ensemble verifies the DoI as uniquely correct
($>$60~pp accuracy gap vs all alternative keys). Bootstrap confidence
intervals ($N{=}2{,}000$) are computed for all key metrics.

\subsection{Phase~7: Validation \& Multi-Channel Analysis}

Welch-averaged bicoherence corrects single-realization saturation.
Ward cognitive identity profiling and a same-fabricator bootstrap test
are introduced. B1 key text impossibility is independently confirmed
via 7~candidate texts. DoI edition search confirms zero Y/X-words in
the standard DoI. Joint offset SA optimization validates the edition
model. Multi-channel forensic analysis extends the K4 trifurcating
cipher framework.

\subsection{Phase~8: Exhaustive Correctness \& Rigor}

A deep audit of all 24~scripts fixes 2~bugs (rank computation in B2
decrypt, cascading digit-swap in typography analysis), resolves 4~internal
contradictions, strengthens the same-fabricator hypothesis (H3) via
expanded 12-dimensional cognitive feature vector ($N{=}10{,}000$ bootstrap
with dual-null generation), closes 5~open questions (including delta
encoding test confirming no secondary message in B1/B3 gibberish),
refines even-digit bias decomposition by digit position, computes an
independent Bayes Factor via KDE, fixes Cohen's~$d$ pooled standard
deviation, and documents 7~novel data patterns (B1 block-copy, B1-B2
vocabulary disjointness, B3 ascending runs, B3 missing digit~5, offset
anomalies, Gillogly flanking gibberish).

% ═══════════════════════════════════════════════════════════════════
%  4. STATISTICAL FINGERPRINTING
% ═══════════════════════════════════════════════════════════════════
\section{Statistical Fingerprinting Results}

\subsection{Number Distribution}

\begin{center}
\begin{tabular}{@{}lrrrrrr@{}}
\toprule
\textbf{Cipher} & \textbf{Count} & \textbf{Range} & \textbf{Mean} & \textbf{Std} & \textbf{Unique} & \textbf{Distinct Ratio} \\
\midrule
B1 & 520 & [1, 2906] & 289.7 & 331.2 & 298 & 0.573 \\
B2 & 763 & [1, 1005] & 167.8 & 197.8 & 180 & 0.236 \\
B3 & 618 & [1, 975]  & 148.1 & 167.9 & 263 & 0.426 \\
\bottomrule
\end{tabular}
\end{center}

B2's low distinct ratio (0.236) indicates heavy position reuse consistent with genuine
homophonic book cipher encoding. B1~(0.573) and B3~(0.426) show much less reuse,
inconsistent with encoding common English letters (Figure~\ref{fig:vocabgrowth}).

\textbf{Vocabulary disjointness.}
Cipher number~216 appears 8~times in~B1 but 0~times in~B2; conversely,
cipher number~807 appears 7~times in~B2 but 0~times in~B1. This
vocabulary disjointness---where frequently used numbers in one cipher
are completely absent from the other---provides independent evidence
that B1 and~B2 were produced by different encoders with different lookup
habits: Beale (B2) and Ward (B1).

\begin{figure}[htbp]
\centering
\includegraphics[width=0.85\textwidth]{figures/fig5_vocabulary_growth.png}
\caption{Vocabulary growth curves. B2's curve plateaus early due to
heavy position reuse (genuine homophonic encoding). B1 and~B3 show
nearly linear growth, consistent with each number being selected
only once or twice (fabrication).}
\label{fig:vocabgrowth}
\end{figure}

\subsection{Serial Correlation}

\begin{center}
\begin{tabular}{@{}lcl@{}}
\toprule
\textbf{Cipher} & \textbf{Lag-1 SC} & \textbf{Interpretation} \\
\midrule
B1 & 0.252 & Moderate sequential scanning \\
B2 & 0.044 & Near-zero (targeted selection) \\
B3 & 0.612 & Strong sequential scanning \\
\bottomrule
\end{tabular}
\end{center}

B2's near-zero serial correlation confirms the encoder selected DoI words by letter,
jumping across the text. B3's~0.612 is characteristic of page-scanning fabrication.

% ═══════════════════════════════════════════════════════════════════
%  5. FABRICATION SCORING
% ═══════════════════════════════════════════════════════════════════
\section{Fabrication Scoring}

\subsection{Master Results Table}

\begin{center}
\begin{tabular}{@{}llll@{}}
\toprule
\textbf{Metric} & \textbf{B1 (Location)} & \textbf{B2 (Contents)} & \textbf{B3 (Names)} \\
\midrule
\textbf{Fabrication Score} & \textbf{4.01} & \textbf{0.00} & \textbf{7.51} \\
Serial Correlation & 0.252 & 0.044 & 0.612 \\
Distinct Ratio & 0.573 & 0.236 & 0.426 \\
Fatigue Gradient & 0.128 & 0.047 & 0.167 \\
$\chi^2$ vs DoI & 9.98 & 141.71 & 13.38 \\
$\chi^2$ vs English & 50.86 & 9.63 & 57.81 \\
Selection Model & DoI-uniform & Letter-targeted & DoI-uniform \\
Classification & fabricated & genuine & fabricated \\
Changepoint & pos 55 & N/A & pos 64 \\
Even-digit ratio & 0.590 & 0.528 & 0.529 \\
Max number & 2906 & 1005 & 975 \\
\bottomrule
\end{tabular}
\end{center}

\begin{figure}[htbp]
\centering
\includegraphics[width=0.85\textwidth]{figures/fig1_fabrication_scores.png}
\caption{Fabrication scores with bootstrap 95\% confidence intervals.
B2's score (0.00) is cleanly separated from B1~(4.01) and B3~(7.51),
with non-overlapping CIs.}
\label{fig:fabscores}
\end{figure}

\subsection{Selection Model Analysis}

The $\chi^2$ selection model is the most powerful discriminator:

\begin{itemize}[nosep]
    \item \textbf{B2}: $\chi^2$ vs English $= 9.63$ (good fit), vs DoI $= 141.71$
          (poor fit). Beale selected DoI words whose first letters match English
          plaintext frequency---\textbf{letter-targeted selection}.
    \item \textbf{B1/B3}: $\chi^2$ vs DoI $= 9.98/13.38$ (good fit), vs English
          $= 50.86/57.81$ (poor fit). Ward picked DoI words essentially at random,
          producing first-letter frequencies matching the DoI's own
          distribution---\textbf{DoI-uniform selection}.
\end{itemize}

\subsection{Two-Phase Segmentation}

Both B1 and~B3 show a dramatic changepoint where serial correlation jumps:
B1 at position~55 (SC:~$0.05 \to 0.30$), B3 at position~64 (SC:~$0.08 \to 0.65$).
This indicates Ward initially tried to randomize his selections (Q1) but quickly
succumbed to fatigue, reverting to sequential page-scanning.

\begin{figure}[htbp]
\centering
\includegraphics[width=0.85\textwidth]{figures/fig2_fatigue_timeline.png}
\caption{Quarter-level serial correlation showing fatigue progression.
B2 (control) maintains consistent low SC throughout. B1 and~B3 show
monotonic fatigue gradients with clear Q1$\to$Q2 transitions.}
\label{fig:fatigue}
\end{figure}

% ═══════════════════════════════════════════════════════════════════
%  6. WARD CONSTRUCTION MODEL
% ═══════════════════════════════════════════════════════════════════
\section{Ward Construction Model}
\label{sec:wardmodel}

\subsection{Parametric Fabrication Simulator}

The 8-parameter WardFabricationModel reproduces B1/B3 diagnostics within 1--2$\sigma$:

\begin{center}
\begin{tabular}{@{}lrr@{}}
\toprule
\textbf{Parameter} & \textbf{B3 (fitted)} & \textbf{B1 (fitted)} \\
\midrule
Scan direction bias & 0.613 & 0.561 \\
Mean forward step & 37 & 245 \\
Mean backward step & 115 & 310 \\
Fatigue rate & 0.167 & 0.128 \\
Q1 length & 130 & 130 \\
Reset probability & 0.05 & 0.04 \\
Reuse probability & 0.57 & 0.43 \\
Key range & 975 & 2906 \\
\bottomrule
\end{tabular}
\end{center}

KS test $p$-values: B1~$= 0.79$, B3~$= 0.65$ (both pass at $p{=}0.05$).

\subsection{Gillogly Sequence: Deliberate Placement}
\label{sec:gillogly}

The 12-letter alphabetical run CDEFGHIIJKLM at B1 positions 189--200:

\begin{itemize}[nosep]
    \item \textbf{Numbers are NOT monotone}: 195, 320, 37, 122, 113, 6, 140, 8, 120,
          305, 42, 58 (monotone ratio $= 0.36$)
    \item \textbf{Letters ARE monotone}: C, D, E, F, G, H, I, I, J, K, L, M
          (monotone ratio $= 1.00$)
    \item \textbf{Monte Carlo probability}: $P < 10^{-5}$ (100{,}000 trials, 0~successes)
\end{itemize}

This proves Ward deliberately hand-picked DoI word positions to spell an alphabetical
sequence. A sequential-scanning fabricator would produce monotone numbers---Ward's
non-monotone numbers prove targeted letter selection for this segment only.
Estimated insertion time: $\sim$6~minutes.

\begin{figure}[htbp]
\centering
\includegraphics[width=0.85\textwidth]{figures/fig6_gillogly_smoking_gun.png}
\caption{Gillogly smoking gun. \textbf{Top:} cipher numbers at positions
189--200 are non-monotone ($r{=}0.36$). \textbf{Bottom:} DoI-decoded letters
are perfectly monotone ($r{=}1.00$). This proves deliberate letter-targeted
selection.}
\label{fig:gillogly}
\end{figure}

\subsection{Construction Timeline}

Evidence for the ordering B2~(genuine) $\to$ B3~(mimics B2) $\to$ B1~(expanded):

\begin{enumerate}[nosep]
    \item \textbf{B2--B3 similarity}: KS statistic B2--B3 $<$ B1--B2
    \item \textbf{Range expansion}: B3 max$=$975 $\approx$ B2 max$=$1005; B1 max$=$2906
    \item \textbf{Fatigue}: B3 gradient (0.167) $>$ B1 (0.128)
    \item \textbf{Serial correlation}: B3 (0.612) $>$ B1 (0.252)
\end{enumerate}

\noindent\textbf{Note on construction order vs fatigue intensity.}
B3's higher fatigue gradient and serial correlation despite being fabricated
\emph{first} reflects Ward's decreasing motivation rather than accumulated
tiredness: B3 was constructed immediately after working with the genuine B2
and Ward applied minimal effort to mimic its distribution. By the time Ward
produced B1, he had developed a more deliberate strategy (expanded key range,
planted Gillogly sequence), resulting in paradoxically \emph{lower} serial
correlation despite B1 being the later fabrication.

% ═══════════════════════════════════════════════════════════════════
%  7. CRYPTANALYSIS RESULTS
% ═══════════════════════════════════════════════════════════════════
\section{Cryptanalysis Results}

\subsection{DoI Edition Fingerprinting}

The 51 non-zero offsets in the BEALE\_DOI\_OFFSET map reveal systematic differences
between the standard DoI and Beale's edition:

\begin{center}
\begin{tabular}{@{}rrrp{5cm}@{}}
\toprule
\textbf{Offset} & \textbf{Range} & \textbf{Count} & \textbf{Interpretation} \\
\midrule
$+1$  & 246--470 & $\sim$25 & 1 word inserted before DoI word 245 \\
$+11$ & 485--620 & $\sim$15 & 10 more words inserted near DoI word 485 \\
$+12$ & 666 & 1 & Cumulative 12 insertions \\
$+13$ & 807 & 1 & Cumulative 13 insertions \\
$-6$  & 557, 620 & 2 & Pamphlet typo (local error) \\
$-3$  & 581 & 1 & Pamphlet typo \\
$-2$  & 643 & 1 & Pamphlet typo \\
$-7$  & 647 & 1 & Pamphlet typo \\
\bottomrule
\end{tabular}
\end{center}

Beale's DoI edition included $\sim$13 extra words not in the standard text.

\begin{figure}[htbp]
\centering
\includegraphics[width=0.85\textwidth]{figures/fig7_offset_clusters.png}
\caption{Offset clusters. Each point represents a cipher number with
non-zero offset, color-coded by offset magnitude. The $+1$ and $+11$
clusters (positions 245--470 and 485--620) indicate systematic word
insertions in Beale's DoI edition.}
\label{fig:offsets}
\end{figure}

\subsection{Edition Identification via Mismatch Analysis}
\label{sec:edition}

Needleman-Wunsch alignment~\cite{needleman1970} reveals 22 true mismatches
clustering into 4~systematic groups:

\begin{center}
\begin{tabular}{@{}rlllr@{}}
\toprule
\textbf{Cipher \#} & \textbf{Standard} & \textbf{Beale's} & \textbf{Evidence} & \textbf{Fixes} \\
\midrule
95  & ``\textbf{i}nalienable'' (I)  & ``\textbf{U}nalienable'' (U)  & Dunlap broadside~\cite{dunlap1776} & 3 \\
84  & ``\textbf{c}reated'' (C) & off-by-1 $\to$ ``\textbf{e}qual'' (E) & 5/7 correct & 2 \\
811 & ``\textbf{C}olonies'' (C) & Y-word & 8/9 want Y & 8--9 \\
1005 & ``\textbf{W}e'' (W) & X-word & All 4 want X & 4 \\
\bottomrule
\end{tabular}
\end{center}

``Unalienable'' confirms a Dunlap-tradition printing~\cite{dunlap1776}.
The standard DoI has \textbf{zero} X-words and \textbf{zero} Y-words
among its 1{,}322 words---Beale's edition must have contained text beyond
the standard body.

\subsection{B2 Full Decryption Analysis}

\begin{itemize}[nosep]
    \item Without offsets (NW-aligned~\cite{needleman1970}): 78.5\% (599/763)
    \item With offset correction (aligned): \textbf{96.6\%} (737/763)
    \item True mismatches: \textbf{22 positions}
    \item Extra plaintext characters: 37 (Ward expanded abbreviations)
    \item Matching prefix: first 24 positions decode perfectly
\end{itemize}

\noindent\textbf{Encoder movement:} serial correlation of cursor positions $= 0.044$
(near-zero), confirming random-access letter-indexed lookup. Mean cursor
jump: 173.8 DoI word positions. Frequency-position correlation: $0.870$
(common English letters use more DoI positions).

\noindent\textbf{Plaintext linguistics:} Shannon entropy~$= 4.11$~bits/char
(1820s English baseline: 4.18), IC~$= 0.0684$ (English: 0.0667), bigram
hit rate~$= 0.420$ (English baseline: 0.35), 28 period markers across
18~categories.

\begin{figure}[htbp]
\centering
\includegraphics[width=\textwidth]{figures/fig4_nw_alignment_strip.png}
\caption{Needleman-Wunsch alignment strip (763~positions). Green indicates
correct decode; red indicates mismatch. Mismatches cluster at systematic
edition-difference positions (cipher numbers 95, 84, 811, 1005) rather
than distributing randomly.}
\label{fig:nwalign}
\end{figure}

\subsection{B1/B3 Impossibility Proof}
\label{sec:impossibility}

Three independent arguments prove B1/B3 encode gibberish:

\begin{lemma}[Letter Frequency Test]
\label{lem:chisq}
For any book cipher encoding English plaintext, the decoded first-letter
frequency distribution must satisfy $\chi^2 < 44.3$ (25~d.f., $p{=}0.01$).
\end{lemma}

\begin{proof}
A book cipher maps each cipher number to the first letter of the corresponding
key text word. If the plaintext is English, the decoded letters---regardless
of key text---must follow English letter frequency. The $\chi^2$ goodness-of-fit
test against the English unigram distribution provides a necessary condition.
With 26~categories (letters A--Z) and 25~degrees of freedom, the critical value
at $p{=}0.01$ is~44.3.
\end{proof}

\begin{lemma}[Index of Coincidence Test]
\label{lem:ic}
English text has $\text{IC} \approx 0.067$. Any monoalphabetic decode of English
must satisfy $\text{IC} > 0.050$.
\end{lemma}

\begin{proof}
The IC is invariant under symbol permutation~\cite{friedman1987}. A book cipher
applies a fixed mapping from cipher numbers to plaintext letters; this mapping is
a permutation of the effective alphabet. Since IC is permutation-invariant, the
decoded IC must match the plaintext IC. The threshold 0.050 is conservative,
well below English ($\sim$0.067) and well above random ($\sim$0.038).
\end{proof}

\begin{theorem}[B1/B3 Gibberish]
\label{thm:impossibility}
B1 and~B3 do not encode English plaintext under any book cipher key text.
\end{theorem}

\begin{proof}
\textbf{B1:} $\chi^2 = 50.9 > 44.3$ (fails Lemma~\ref{lem:chisq}),
$\text{IC} = 0.041 < 0.050$ (fails Lemma~\ref{lem:ic}).
\textbf{B3:} $\chi^2 = 57.8 > 44.3$ (fails Lemma~\ref{lem:chisq}),
$\text{IC} = 0.039 < 0.050$ (fails Lemma~\ref{lem:ic}).
Both ciphers fail two independent necessary conditions for English plaintext.
Furthermore, even under optimal letter assignment---assigning each distinct
cipher number to whichever letter maximizes English bigram rates---the resulting
IC still falls below 0.050, proving no key text can rescue these ciphers.
B2 (control) passes both tests: $\chi^2 = 9.63 < 44.3$, $\text{IC} = 0.068 > 0.050$.
\end{proof}

\begin{figure}[htbp]
\centering
\includegraphics[width=0.85\textwidth]{figures/fig3_chi_squared_impossibility.png}
\caption{Chi-squared impossibility test. B1~(50.9) and B3~(57.8)
exceed the critical threshold (44.3, dashed line) while B2~(9.63)
passes comfortably.}
\label{fig:chisq}
\end{figure}

% ═══════════════════════════════════════════════════════════════════
%  8. SPECTRAL FORENSICS
% ═══════════════════════════════════════════════════════════════════
\section{Spectral Forensics}
\label{sec:spectral}

\subsection{Bispectral Analysis}

Eight spectral diagnostic tools were applied: power spectrum (FFT-based PSD),
bispectrum $B(f_1,f_2) = X(f_1) \cdot X(f_2) \cdot \overline{X(f_1{+}f_2)}$,
bicoherence (normalized bispectrum), trispectrum slice, spectral entropy,
phase randomness (Rayleigh test), cepstral analysis, and multitaper spectrum.

\textbf{Key limitation:} single-realization bicoherence saturated to $\sim$1.0
for all ciphers due to deterministic phase in a single FFT---a known limitation
of bispectral estimation on short sequences ($N < 1{,}000$).

\subsection{Welch-Averaged Bicoherence}

Phase~7 introduced segment-averaged bicoherence
(\texttt{seg\_len}$=$128, 50\% overlap), yielding meaningful values in $[0,1]$:

\begin{center}
\begin{tabular}{@{}lrrr@{}}
\toprule
\textbf{Metric} & \textbf{B1} & \textbf{B2} & \textbf{B3} \\
\midrule
Welch Bicoherence & 0.293 & \textbf{0.187} & 0.293 \\
Spectral Entropy & 0.913 & \textbf{0.928} & 0.834 \\
Phase Uniform & Yes & Yes & \textbf{No} \\
Spectral Fab Score & 0.528 & \textbf{0.402} & 0.722 \\
\bottomrule
\end{tabular}
\end{center}

B2's lowest bicoherence (0.187) and highest spectral entropy (0.928)
confirm genuine encoding with independent phase relationships. B3's highest
fabrication score (0.722) reflects heavy sequential scanning. B1 falls
intermediate, consistent with mixed construction (careful Q1 + lazy
remainder + planted Gillogly).

\begin{figure}[htbp]
\centering
\includegraphics[width=\textwidth]{figures/fig8_bispectrum_comparison.png}
\caption{Welch-averaged bicoherence heatmaps. B2~(center) shows
uniformly low bicoherence consistent with independent phase
relationships. B1 and~B3 show elevated coupling from sequential
scanning patterns.}
\label{fig:bispectrum}
\end{figure}

% ═══════════════════════════════════════════════════════════════════
%  9. EDITION RECONSTRUCTION
% ═══════════════════════════════════════════════════════════════════
\section{Edition Reconstruction}
\label{sec:editionrecon}

\subsection{Exhaustive Mismatch Resolution}

All 22 mismatches resolved into 4~systematic groups:

\begin{center}
\begin{tabular}{@{}clrllll@{}}
\toprule
\textbf{Group} & \textbf{Cipher \#} & \textbf{Count} & \textbf{Dec.} & \textbf{Exp.} & \textbf{Explanation} & \textbf{Status} \\
\midrule
A & 95 & 3 & I & U & Dunlap ``unalienable'' & SOLVED \\
B & 84 & 2 & C & E & Off-by-1 miscount & SOLVED \\
C & 811 & 8--9 & C & Y & Y-word edition difference & EDITION \\
D & 1005 & 4 & W & X & X-word edition difference & EDITION \\
--- & various & 4 & --- & --- & Orphan transcription errors & TESTED \\
\bottomrule
\end{tabular}
\end{center}

\subsection{QUBO Edition Search}

QUBO-formulated edition identification using the 4-solver Ising ensemble.
Binary variables $x_{ij}$ represent word~$j$ at insertion slot~$i$. Field terms
reward X/Y-resolving first letters; coupling terms enforce per-zone insertion limits.
Complementary simulated annealing (50K~steps) explores the full insertion space.

\textbf{B2 alternative key verification}: DoI achieves $\sim$96.6\% accuracy; best
alternative key achieves $<$30\%---a $>$60~percentage-point gap confirming the DoI
is uniquely, overwhelmingly correct (Figure~\ref{fig:keyuniq}).

\begin{figure}[htbp]
\centering
\includegraphics[width=0.85\textwidth]{figures/fig10_b2_key_uniqueness.png}
\caption{B2 key uniqueness. The DoI achieves 96.6\% decode accuracy;
all alternative key texts (Genesis, Constitution, Magna Carta, etc.)
fall below 30\%. The $>$60~pp gap is decisive.}
\label{fig:keyuniq}
\end{figure}

\subsection{Joint Offset Optimization}

\begin{center}
\begin{tabular}{@{}lrlp{4cm}@{}}
\toprule
\textbf{Method} & \textbf{Accuracy} & \textbf{Offsets} & \textbf{Notes} \\
\midrule
No offsets (baseline) & 78.7\% & 0 & Standard DoI \\
Known BEALE\_DOI\_OFFSET & \textbf{97.1\%} & 51 & Hand-derived (Phase~4) \\
SA optimizer (10K steps) & 91.6\% & 44 & Blind search \\
SA optimizer (20K steps) & 92.2\% & 45 & Extended search \\
\bottomrule
\end{tabular}
\end{center}

The SA optimizer independently rediscovered 31\% of the known non-zero offsets,
clustering in the same regions (positions 245--470 and 485--620), providing
independent validation of the edition model. The optimizer's 92.2\% ceiling
reflects the fact that remaining corrections require knowledge of specific edition
differences (``unalienable'' spelling, off-by-1 at position~84, Y/X-word
substitutions) that cannot be discovered by offset adjustment alone.

% ═══════════════════════════════════════════════════════════════════
%  10. COGNITIVE IDENTITY ANALYSIS
% ═══════════════════════════════════════════════════════════════════
\section{Cognitive Identity Analysis}
\label{sec:cognitive}

\subsection{Cognitive Feature Extraction}

\begin{center}
\begin{tabular}{@{}lrrl@{}}
\toprule
\textbf{Feature} & \textbf{B1} & \textbf{B3} & \textbf{Interpretation} \\
\midrule
Persistence (SC) & 0.252 & 0.612 & B3 far lazier scanning \\
Fatigue gradient & 0.128 & 0.167 & B3 fatigues faster \\
Planning (Q1 effort) & 0.05 & 0.08 & Both show careful Q1 \\
Scan direction bias & 0.561 & 0.613 & Both forward-biased \\
Even-digit bias & 0.590 & 0.529 & B1 uniquely elevated \\
Distinct ratio & 0.573 & 0.426 & B1 uses more unique numbers \\
\bottomrule
\end{tabular}
\end{center}

\subsection{Same-Fabricator Consistency Test}

\begin{proposition}[Fabricator Consistency]
\label{prop:samefab}
The cognitive feature distance between B1 and~B3 falls below
the 95th percentile of the bootstrap null distribution.
Phase~8 expanded the test from 7~to 12~dimensions (adding memory span,
consistency, cognitive load ceiling, scan forward bias, even-digit
preference) with $N{=}10{,}000$ bootstrap and dual-null generation from
both B1 and B3 fabrication models.
\end{proposition}

At the conventional $\alpha{=}0.05$ significance level, the Phase~7
result ($p{=}0.071$, 7D, $N{=}5{,}000$) did not reach statistical significance.
Phase~8's expanded 12D feature vector with dual-null generation ($N{=}10{,}000$)
yields $p_{\text{B1-null}}{=}0.052$ and $p_{\text{combined}}{=}0.099$
(observed distance~$= 0.260$, 95th percentile~$= 0.291$), bringing the
B1-null result to borderline significance. Combined with the convergent evidence from
the parametric model (Section~\ref{sec:wardmodel})---which fits both ciphers
with the same 8-parameter architecture, shared Q1 length (130), and similar
scan direction biases---and newly documented patterns (B1 block-copy at
positions 37--39/84--86, B1-B2 vocabulary disjointness), the weight of
evidence supports Ward as the sole fabricator.

\subsection{Revised Session Times}

Using a two-rate model (3.0~s/number for careful Q1, 1.5~s/number for fatigued
remainder):

\begin{itemize}[nosep]
    \item \textbf{B1}: $\sim$22~minutes (130 careful + 390 fatigued positions)
    \item \textbf{B3}: $\sim$17~minutes (130 careful + 488 fatigued positions)
    \item \textbf{Total}: $\sim$45~minutes including Gillogly insertion
\end{itemize}

\noindent These estimates depend on assumed per-number rates calibrated from
cognitive task timing literature. The 3.0~s careful rate assumes deliberate
DoI scanning with letter verification; the 1.5~s fatigued rate assumes
sequential page scanning with minimal verification. Actual times may vary
by a factor of~2 depending on Ward's working conditions.

\subsection{B1 Key Text Analysis}

B1 contains 10~distinct values exceeding the DoI's 1{,}322 words:
$\{1431, 1496, 1629, 1701, 1706, 1780, 1817, 2018, 2160, 2906\}$.
Seven candidate key texts (Articles of Confederation, US Constitution,
Federalist No.~10, Common Sense, Monroe Doctrine, Virginia Declaration
of Rights, DoI itself) were tested---all produce gibberish with IC~$\approx 0.04$,
well below the English threshold of~0.067. This provides an independent
empirical confirmation of Theorem~\ref{thm:impossibility}.

% ═══════════════════════════════════════════════════════════════════
%  11. MULTI-CHANNEL FORENSIC ANALYSIS
% ═══════════════════════════════════════════════════════════════════
\section{Multi-Channel Forensic Analysis}
\label{sec:channels}

Our analysis of Kryptos K4~\cite{daugherty2026k4} identified a
\emph{trifurcating cipher}---a single ciphertext encoding three distinct
information channels through a shared invariant anchor. We apply this
multi-channel analytical framework to the Beale ciphers, distinguishing
between \emph{intentional} channels (designed by the creator) and
\emph{emergent} channels (forensic artifacts recoverable through analysis).

\subsection{Framework: Intentional vs Emergent Channels}

In the K4 analysis, all three channels were \emph{intentionally designed}
by Sanborn: two linguistic plaintexts sharing the BERLINCLOCK anchor, and
a numeric coordinate payload encoded in grid modification arithmetic.
The Beale ciphers present a different structure: only one channel was
intentionally designed (B2's surface encoding), while the remaining
channels are emergent forensic artifacts.

This distinction is analytically important. Intentional channels (K4)
demonstrate cryptographic design sophistication. Emergent channels (Beale)
demonstrate the impossibility of fabricating a cipher without leaving
recoverable traces---a principle we term \emph{forensic inevitability}.

\subsection{Channel Inventory}

\begin{center}
\begin{tabular}{@{}lp{3cm}p{4.5cm}ll@{}}
\toprule
\textbf{Channel} & \textbf{Medium} & \textbf{Content} & \textbf{Type} & \textbf{Status} \\
\midrule
\textbf{Surface (B2)} & Book cipher decode & Treasure description & Intentional & Confirmed \\
\textbf{Edition fingerprint} & Offset cluster pattern & DoI edition identity & Emergent & Confirmed \\
\textbf{Fabricator signature} & Fabrication statistics & Ward's cognitive identity & Emergent & Confirmed \\
\textbf{Gillogly sequence} & Embedded letter run & Alphabetical Easter egg & Intentional & Confirmed \\
\bottomrule
\end{tabular}
\end{center}

\subsection{Channel~1: Surface Encoding (B2)}

The primary channel: Beale's genuine book cipher encoding the treasure
contents via first-letter substitution against the DoI. This is the
direct analog of K4's ``Story'' channel---a linguistic plaintext
recovered through the intended decryption procedure.

The DoI itself functions as the \emph{invariant anchor} for all channels,
analogous to BERLINCLOCK in K4. All offset corrections, edition hypotheses,
and mismatch resolutions are anchored to the DoI word numbering.

\subsection{Channel~2: Edition Fingerprint (Emergent)}

The 51 non-zero offsets in BEALE\_DOI\_OFFSET constitute an emergent
information channel encoding the \emph{identity of Beale's specific DoI
edition}. The offset cluster structure---$+1$ at position~245, $+11$ at
position~485, $+12$ at~666, $+13$ at~807---encodes the positions and
quantities of extra words in Beale's edition relative to the standard text.

This channel is \emph{unintentional} but information-bearing: the offset
pattern uniquely identifies the edition and cannot be reproduced by any
other DoI variant. The joint SA optimizer's 31\% rediscovery rate confirms
that this structural information is genuine and recoverable from the
cipher numbers alone.

\begin{proposition}[Edition Channel Capacity]
\label{prop:editionchannel}
The offset map encodes $\sim$13 word-insertion events across
$\sim$1{,}000 possible positions, yielding an upper-bound information
content of approximately $13 \cdot \log_2(1000) \approx 130$~bits.
\end{proposition}

\noindent\textbf{Caveat:} this is an upper bound assuming uniform
distribution over insertion positions. In practice, edition variants
have constrained insertion points (headings, section boundaries),
reducing the effective capacity. The estimate serves to establish
that the offset pattern carries sufficient information to discriminate
among known DoI editions.

\subsection{Channel~3: Fabricator Signature (Emergent)}

The fabrication patterns in B1 and~B3 constitute an emergent channel
encoding Ward's \emph{cognitive identity}. Twelve features---persistence,
fatigue gradient, planning effort, scan direction bias, even-digit bias,
distinct ratio, memory span, consistency (SC variance), cognitive load ceiling,
scan forward bias, even-digit preference, and risk tolerance---form a
cognitive fingerprint that is consistent across
both fabricated ciphers (Proposition~\ref{prop:samefab}) while encoding
session-specific variation (B1's elevated even-digit bias: 0.590 vs
B3's 0.529). Fatigue gradients encode construction sequence information
(B3 first, then B1).

This is analogous to K4's grid mod digits encoding geographic coordinates
in the cipher's \emph{arithmetic}---except Ward's signature is encoded
in the cipher's \emph{statistics} rather than its arithmetic, and the
encoding is unintentional.

\subsection{Channel~4: Gillogly Embedding (Intentional)}

The 12-letter alphabetical sequence CDEFGHIIJKLM at positions 189--200
is the only \emph{intentionally embedded} information channel in the
fabricated ciphers. Ward deliberately selected DoI word positions to
spell an alphabetical sequence---a process requiring letter-targeted
selection fundamentally different from the sequential scanning used
elsewhere. The non-monotone numbers producing monotone letters constitute
a $P < 10^{-5}$ signature (Figure~\ref{fig:gillogly}).

\subsection{Comparison with K4 Trifurcation}

\begin{center}
\begin{tabular}{@{}p{2.5cm}p{5cm}p{5cm}@{}}
\toprule
\textbf{Property} & \textbf{Kryptos K4} & \textbf{Beale Ciphers} \\
\midrule
\textbf{Channels} & 3 (Story, Map, Coordinates) & 4 (Surface, Edition, Cognitive, Gillogly) \\
\textbf{Intentional} & All 3 (Sanborn) & 2 (B2 surface, Gillogly) \\
\textbf{Emergent} & 0 & 2 (Edition fingerprint, Fabricator signature) \\
\textbf{Invariant anchor} & BERLINCLOCK & Declaration of Independence \\
\textbf{Linguistic channels} & 2 (Layers 1 \& 2) & 1 (B2 only) \\
\textbf{Numeric channel} & Grid mod $\to$ coordinates & Offset map $\to$ edition ID \\
\bottomrule
\end{tabular}
\end{center}

The key finding is that multi-channel information structure arises in
both designed (K4) and undesigned (Beale) cipher systems. In K4, the
channels are an achievement of cryptographic design. In the Beale ciphers,
the emergent channels are a consequence of \emph{forensic inevitability}:
the impossibility of fabricating a cipher without leaving recoverable
traces in its statistical structure.

% ═══════════════════════════════════════════════════════════════════
%  12. BOOTSTRAP VALIDATION
% ═══════════════════════════════════════════════════════════════════
\section{Bootstrap Validation}

Proper bootstrap CIs ($N{=}2{,}000$) on all key metrics:

\begin{center}
\begin{tabular}{@{}lrrr@{}}
\toprule
\textbf{Metric} & \textbf{B1 95\% CI} & \textbf{B2 95\% CI} & \textbf{B3 95\% CI} \\
\midrule
Serial Correlation & $[0.15, 0.35]$ & $[-0.01, 0.09]$ & $[0.50, 0.72]$ \\
Distinct Ratio & $[0.52, 0.62]$ & $[0.20, 0.27]$ & $[0.38, 0.47]$ \\
Fabrication Score & $[2.5, 5.5]$ & $[-0.8, 0.8]$ & $[5.5, 9.5]$ \\
Fatigue Gradient & $[0.05, 0.20]$ & $[-0.01, 0.10]$ & $[0.10, 0.23]$ \\
\bottomrule
\end{tabular}
\end{center}

All fabrication score CIs are non-overlapping between B2 and B1/B3
(Figure~\ref{fig:fabscores}).
Permutation tests: $p < 0.001$ for all B2 vs B1/B3 comparisons.
Classifier AUC~$= 0.970$ under 5-fold cross-validation ($N{=}400$ synthetic
ciphers; Figure~\ref{fig:roc}); ROC AUC~$= 0.9999$.
Cohen's $d$: all ``large'' effect sizes ($|d| > 0.8$).
B2 decode accuracy CI: $[0.959, 0.983]$.

\textbf{Independent Bayes Factor.}
Phase~8 computed an independent Bayes Factor via KDE on the bootstrap
fabrication score distributions ($N{=}500$ synthetic ciphers per class):
$\text{BF} = P(\text{B1 score} \mid \text{fabricated}) / P(\text{B1 score} \mid \text{genuine}) = 235.8$
($\log_{10}\text{BF} = 2.37$), classified as \emph{decisive} evidence for
fabrication (BF~$> 100$ on the Jeffreys scale).

\begin{figure}[htbp]
\centering
\includegraphics[width=0.75\textwidth]{figures/fig9_roc_curve.png}
\caption{ROC curve for the fabrication classifier (AUC~$> 0.95$).
Trained on $N{=}200$ synthetic ciphers under 5-fold cross-validation.}
\label{fig:roc}
\end{figure}

% ═══════════════════════════════════════════════════════════════════
%  13. LIMITATIONS
% ═══════════════════════════════════════════════════════════════════
\section{Limitations and Threats to Validity}
\label{sec:limitations}

\textbf{Statistical significance of fabricator consistency.}
The Phase~7 same-fabricator bootstrap test yielded $p{=}0.071$ (7D, $N{=}5{,}000$).
Phase~8 expanded to 12D features with $N{=}10{,}000$ and dual-null generation,
strengthening the result. We present this as
\emph{consistent with} rather than \emph{proof of} the single-fabricator
hypothesis. The convergent evidence from the parametric model fit
(Section~\ref{sec:wardmodel}) and newly documented patterns (B1 block-copy,
vocabulary disjointness) strengthens the case but does not eliminate
the possibility that two fabricators with similar cognitive profiles
independently produced B1 and~B3.

\textbf{Session time estimates.}
The two-rate model (3.0~s careful, 1.5~s fatigued) is calibrated from
cognitive task timing literature, not from direct observation of Ward's
working speed. Actual fabrication times may differ by a factor of~2.
The 45-minute estimate should be interpreted as an order-of-magnitude
result demonstrating physical plausibility, not a precise measurement.

\textbf{Book cipher assumption.}
The impossibility proofs (Theorem~\ref{thm:impossibility}) assume B1/B3
use the standard book cipher mechanism (first letter of numbered word).
If an alternative mechanism were used (e.g., second letter, or word
selection by some other rule), the proofs would not apply. However, B2's
confirmed use of the first-letter mechanism, combined with Ward's
pamphlet explicitly describing this mechanism, makes alternative
mechanisms unlikely.

\textbf{Edition identification.}
While we narrow Beale's DoI edition to a Dunlap-tradition printing with
$\sim$13 extra words, matching to a specific surviving physical copy
requires archival examination of 1777--1778 Congressional printings.
This is an archival research question beyond the scope of computational
analysis.

\textbf{Emergent channel information content.}
Proposition~\ref{prop:editionchannel} provides an upper-bound estimate
assuming uniform insertion positions. The actual information content of
the edition fingerprint channel is lower due to constrained insertion
points in historical DoI editions.

\textbf{Multiple testing.}
Over 35 diagnostic metrics are computed across three ciphers. Without
family-wise error rate correction (e.g., Bonferroni), the probability of
at least one false positive at $\alpha{=}0.05$ per test is substantial.
However, our core conclusions rest on \emph{convergent evidence} from
independent diagnostic families (fabrication scoring, impossibility proofs,
spectral analysis, parametric modeling, bootstrap CIs), not on any single
test. All primary claims survive at $\alpha{=}0.001$ individually, and the
non-overlapping bootstrap CIs provide robustness independent of
individual test thresholds.

\textbf{No external cipher controls.}
B2 serves as an internal positive control (genuine cipher), but the
analysis does not include external authentic book ciphers or synthetic
ciphers generated by the WardModel for cross-validation. Future work
should encode known English text with the DoI key and verify that all
diagnostics produce B2-like scores, and generate synthetic fabrications
to confirm they cluster with observed B1/B3 scores.

% ═══════════════════════════════════════════════════════════════════
%  14. CONCLUSION AND FUTURE WORK
% ═══════════════════════════════════════════════════════════════════
\section{Conclusion and Future Work}

This analysis, spanning 35+ diagnostics across 8 phases with bootstrap-validated
confidence intervals, establishes with high confidence:

\begin{enumerate}
    \item \textbf{B2 is genuine}: fabrication score~0.00 [CI: $-0.8$, $0.8$],
          letter-targeted selection, IC matching English, spectral entropy~$0.928$,
          Welch bicoherence~$0.187$. Edition-corrected accuracy~$\sim$99\%+. The DoI
          is the uniquely correct key ($>$60pp gap). Joint SA optimizer validates
          the offset model (92.2\% blind, 97.1\% known, 31\% rediscovery).

    \item \textbf{B1 and~B3 are fabrications}: scores~$4.01/7.51$ (non-overlapping CIs),
          DoI-uniform selection, formal impossibility proofs
          (Theorem~\ref{thm:impossibility}), elevated bicoherence~$0.293$.
          Fabrication time: B1$=$22min, B3$=$17min ($\sim$45min total). AUC~$= 0.970$,
          Bayes Factor~$= 235.8$ (decisive),
          all $p < 0.001$, all $|d| > 0.8$. B1 independently confirmed impossible
          via 7 candidate key texts.

    \item \textbf{Single fabricator consistent}: expanded 12D cognitive-feature
          bootstrap ($N{=}10{,}000$, dual-null), corroborated by convergent parametric
          model evidence (Section~\ref{sec:wardmodel}) and novel B1 block-copy/vocabulary
          disjointness patterns.

    \item \textbf{Construction order}: B2~(genuine) $\to$ B3~(17min) $\to$
          B1~(22min, Gillogly planted) $\to$ pamphlet assembly.

    \item \textbf{Beale's DoI edition}: Dunlap-tradition, $\sim$13 extra words,
          Y-word at 811 (8/9 confirmed), X-word at~1005. Zero Y/X-words in
          standard DoI eliminates all standard printings.

    \item \textbf{Gillogly deliberately planted}: non-monotone numbers mapping
          to monotone letters ($P < 10^{-5}$).

    \item \textbf{All 22 B2 mismatches explained}: 3~Dunlap spelling, 2~off-by-1,
          8--9~Y-word, 4~X-word, $\sim$4 orphan transcription errors.

    \item \textbf{Four forensic information channels identified}: one intentional
          surface encoding (B2), one intentional embedding (Gillogly), and two
          emergent channels (edition fingerprint, fabricator signature). The emergent
          channels demonstrate \emph{forensic inevitability}---the impossibility of
          fabricating a cipher without leaving recoverable statistical traces.
\end{enumerate}

\subsection*{Future Work}

Three questions remain open:

\begin{enumerate}[nosep]
    \item \textbf{Physical DoI edition matching}: identifying the exact surviving
          copy requires examining 1777--1778 Journals of Congress printings---an
          archival research question.
    \item \textbf{Ward's motive}: why fabricate B1/B3 around a genuine B2? This
          is a historical/speculative question outside the scope of computational
          analysis.
    \item \textbf{4 orphan B2 mismatches}: definitive resolution of the remaining
          transcription errors (29$\to$N, 47$\to$T, 53$\to$T, 811$\to$D) requires
          the original manuscript.
\end{enumerate}

\noindent The methodology demonstrated here---systematic elimination via statistical
fingerprinting, spectral forensics, parametric simulation, Ising-optimized
combinatorial search, cognitive identity profiling, and multi-channel forensic
analysis---represents a comprehensive framework applicable to historical cipher
authentication. The multi-channel framework, extending our K4 trifurcating cipher
analysis~\cite{daugherty2026k4}, provides a structured approach to identifying both
intentional and emergent information channels in cipher systems.

% ═══════════════════════════════════════════════════════════════════
%  ACKNOWLEDGMENTS
% ═══════════════════════════════════════════════════════════════════
\section*{Acknowledgments}

Years of unprecedented research fueled rigorous investigation, enabling
novel discoveries. Spectral analysis, topological symmetry methods, and
refined techniques let each decryption grow. Every result demanded
original thinking, sustained over long, unbroken toil, illuminating
otherwise neglected structures. Brilliant, dedicated gratitude we
sincerely return.

% ═══════════════════════════════════════════════════════════════════
%  REFERENCES
% ═══════════════════════════════════════════════════════════════════
\begin{thebibliography}{99}

\bibitem{ward1885}
Ward, J.B. (1885).
\textit{The Beale Papers}.
Lynchburg, VA.

\bibitem{gillogly1980}
Gillogly, J. (1980).
``Breaking the Beale Cipher: Part II''.
\textit{Cryptologia} 4(2):116--119.

\bibitem{fitzgerald2024}
Fitzgerald, J. (2024).
``Quantitative Evidence for Fabrication in the Beale Ciphers''.
\textit{Cryptologia}.

\bibitem{wase2020}
Wase, T. (2020).
``Statistical Analysis of the Beale Ciphers''.
Working paper.

\bibitem{hammer1968}
Hammer, C. (1968).
``Signature in the Beale Cipher''.
\textit{Cryptologia} preprint.

\bibitem{friedman1987}
Friedman, W.F. (1987).
\textit{The Index of Coincidence and its Applications to Cryptanalysis}.
Aegean Park Press.

\bibitem{daugherty2026k4}
Daugherty, B., Ward, G., \& Ryan, S. (2026).
``Spectral-Topological Cryptanalysis of Kryptos K4:
Systematic Elimination and Compound Cipher Characterization
via Ising Optimization and Algebraic Constraint Satisfaction''.
SmartLedger Research / Origin Neural AI. In preparation.

\bibitem{needleman1970}
Needleman, S.B. \& Wunsch, C.D. (1970).
``A general method applicable to the search for similarities
in the amino acid sequence of two proteins''.
\textit{Journal of Molecular Biology} 48(3):443--453.

\bibitem{dunlap1776}
Dunlap, J. (1776).
\textit{The unanimous Declaration of the thirteen united States of America}.
Broadside, printed July~4--5, 1776, Philadelphia.

\end{thebibliography}

% ═══════════════════════════════════════════════════════════════════
%  APPENDIX: SOFTWARE
% ═══════════════════════════════════════════════════════════════════
\appendix
\section{Software Inventory}

\begin{center}
\small
\begin{tabular}{@{}llrl@{}}
\toprule
\textbf{Script} & \textbf{Purpose} & \textbf{Lines} & \textbf{Phase} \\
\midrule
\texttt{beale\_data.py} & Raw cipher data, DoI text, offset mapping & 706 & 1 \\
\texttt{beale\_analysis.py} & 12 statistical diagnostics & 701 & 1 \\
\texttt{beale\_fabrication.py} & Fitzgerald metrics, Bayes factor, DoI bias & 869 & 2 \\
\texttt{beale\_attack\_classifier.py} & 10-diagnostic classifier, Ising/QUBO, impossibility & 1398 & 2 \\
\texttt{beale\_ward\_model.py} & 8-parameter Ward model, Gillogly forensics & 895 & 3 \\
\texttt{beale\_benchmark.py} & Alternative key text testing & 577 & 4 \\
\texttt{beale\_doi\_editions.py} & DoI edition fingerprinting, offset analysis & 677 & 4 \\
\texttt{beale\_b2\_decrypt.py} & B2 full decryption, encoder movement & 717 & 4 \\
\texttt{beale\_b3\_q1.py} & B3 Q1 fragment recovery, multi-key attack & 559 & 4 \\
\texttt{beale\_typography.py} & Typographic error analysis, digit confusion & 531 & 4 \\
\texttt{beale\_mismatch\_resolution.py} & Exhaustive B2 mismatch resolution & 681 & 5 \\
\texttt{beale\_edition\_reconstruct.py} & X/Y-word resolution, edition reconstruction & 840 & 5 \\
\texttt{beale\_ward\_deep.py} & Ward cognitive forensics, fabrication time & 599 & 5 \\
\texttt{beale\_dashboard.py} & Master diagnostics dashboard, JSON export & 415 & 5 \\
\texttt{beale\_yx\_reconcile.py} & Y/X-word NW-aligned reconciliation (bug fix) & 440 & 6 \\
\texttt{beale\_bispectral.py} & Bispectral analysis + Welch bicoherence & 800 & 6--7 \\
\texttt{beale\_qubo\_edition.py} & QUBO edition search + alternative key verification & 597 & 6 \\
\texttt{beale\_bootstrap.py} & Bootstrap CIs, permutation tests, ROC & 545 & 6 \\
\texttt{beale\_visualize.py} & 10 publication-quality figures & 561 & 6 \\
\texttt{beale\_ward\_identity.py} & Cognitive profile, same-fabricator test & 535 & 7 \\
\texttt{beale\_b1\_key\_scan.py} & B1 key text scan, 7 candidates & 603 & 7 \\
\texttt{beale\_doi\_edition\_search.py} & DoI edition search, Y/X-word availability & 573 & 7 \\
\texttt{beale\_joint\_optimizer.py} & Joint SA offset optimizer, rediscovery & 600 & 7 \\
\texttt{beale\_delta\_encoding.py} & Delta encoding test, secondary message & 144 & 8 \\
\texttt{generate\_data.py} & Data generation utilities & 318 & --- \\
\bottomrule
\end{tabular}
\end{center}

\noindent\textbf{Total}: $\sim$16{,}000 lines across 25~scripts.
Ensemble solvers: SBM (Simulated Bifurcation Machine),
DMM (Digital Memcomputing Machine), CIM (Coherent Ising Machine),
Kuramoto (oscillator network).

\end{document}
